% Part: second-order-logic
% Chapter: syntax-and-semantics
% Section: terms-formulas

\documentclass[../../../include/open-logic-section]{subfiles}

\begin{document}

\olfileid{sol}{syn}{frm}

\olsection{项与\printtoken{P}{formula}}

与一阶逻辑类似，二阶逻辑的表达式也是由包含\emph{变元}、\emph{常项符号}、\emph{谓词符号}（有时还有\emph{函数符号}）的基本词汇表构造而成的。由它们，连同逻辑联结词、量词以及括号、逗号等标点符号，即可构成\emph{项}和\emph{公式}。区别在于，除了个体变元外，二阶逻辑还包含关系变元和函数变元，并允许对它们进行量化。因此，二阶逻辑的逻辑符号即一阶逻辑的逻辑符号，再加上：

\begin{enumerate}
\item 对每个元数~$n$，都有一个可数无穷的二阶关系变元集合：$\Obj V_0^n$，$\Obj V_1^n$，$\Obj V_2^n$，\dots
\item 对每个元数~$n$，都有一个可数无穷的二阶函数变元集合：$\Obj u_0^n$，$\Obj u_1^n$，$\Obj u_2^n$，\dots
\end{enumerate}

正如我们用 $x$、$y$、$z$ 作为一阶变元 $\Obj v_i$ 的元变元，我们将用 $X$、$Y$、$Z$ 等作为 $\Obj V_i^n$ 的元变元，用 $u$、$v$ 等作为~$\Obj u_i^n$ 的元变元。

\begin{explain}
二阶语言的\emph{非逻辑符号}的指定方式与一阶语言相同：通过列出其常项符号、函数符号和谓词符号。

在一阶逻辑中，通常包含等词~$\eq$。语言~$\Lang{L}$ 的非逻辑符号对于表达任何有趣的内容至关重要。当然，也存在不使用任何非逻辑符号的语句，但仅仅用~$\eq$ 很难表达任何有趣的内容。在二阶逻辑中，由于我们拥有无穷多的关系变元和函数变元，即使没有专门提供非逻辑符号，我们也能表达一阶语言所能表达的任何内容。
\end{explain}

\begin{defn}[二阶项]
语言~$\Lang L$ 的\emph{二阶项}集合 $\TrmSOL[L]$ 通过在 \olref[fol][syn][frm]{defn:terms} 的基础上添加以下条款定义：
\begin{enumerate}
\item 如果 $u$ 是 $n$ 元函数变元，且 $t_1$、……、$t_n$ 是项，则 $\Atom{u}{t_1, \ldots, t_n}$ 是项。
\end{enumerate}
\end{defn}

\begin{explain}
因此，二阶项与一阶项看起来完全相同，只不过在一阶项包含函数符号~$\Obj{f^n_i}$ 之处，二阶项可代之以函数变元~$\Obj{u^n_i}$。
\end{explain}

\begin{defn}[二阶\usetoken{s}{formula}]
语言~$\Lang L$ 的\emph{二阶公式}集合~$\FrmSOL[L]$ 通过在 \olref[fol][syn][frm]{defn:formulas} 的基础上添加以下条款定义：
\begin{enumerate}
\item 如果 $X$ 是 $n$ 元谓词变元，且 $t_1$、……、$t_n$ 是~$\Lang L$ 的二阶项，则 $\Atom{X}{t_1,\ldots, t_n}$ 是原子公式。

\tagitem{prvAll}{如果 $!A$ 是公式且 $u$ 是函数变元，则 $\lforall[u][!A]$ 是公式。}{}

\tagitem{prvAll}{如果 $!A$ 是公式且 $X$ 是谓词变元，则 $\lforall[X][!A]$ 是公式。}{}

\tagitem{prvEx}{如果 $!A$ 是公式且 $u$ 是函数变元，则 $\lexists[u][!A]$ 是公式。}{}

\tagitem{prvEx}{如果 $!A$ 是公式且 $X$ 是谓词变元，则 $\lexists[X][!A]$ 是公式。}{}
\end{enumerate}
\end{defn}

\end{document}